\section{População e Amostragem}

\begin{itemize}
  \item \textbf{População-alvo:} desenvolvedores, engenheiros (backend/frontend/fullstack/mobile) e arquitetos de software em atividade remunerada no Brasil (empregados, servidores, PJ e autônomos).
  \item \textbf{Inclusão:} (1) atuação remunerada em desenvolvimento de software; (2) uso de ao menos um assistente de IA para código nos últimos 6 meses (habitual ou esporádico; a frequência é covariável, ITEM\_07); (3) atuação no Brasil; (4) idade $\ge 18$ anos (declarada no consentimento).
  \item \textbf{Exclusão:} estudantes sem vínculo profissional; colegas da mesma turma da disciplina; respostas que falhem nos controles de qualidade da Seção~5.1. Os ITEM\_01--04 usam \textit{skip logic}: quem não é elegível encerra o formulário e nada é gravado.
  \item \textbf{Origem dos contatos e estratégia:} amostragem não probabilística, por conveniência com bola de neve: LinkedIn dos integrantes e grupos técnicos de Telegram/WhatsApp; cada respondente é convidado a repassar a pares. Não permite generalização estatística; os resultados descrevem a amostra obtida.
  \item \textbf{Tamanho pretendido:} 60 respostas válidas (mínimo aceitável: 50), com meta de $\ge 80$ respostas brutas para absorver descartes.
  \item \textbf{Margem de erro (95\%, $z = 1{,}96$):} $e = 1{,}96 \times \sqrt{p(1 - p)/n}$. No pior caso ($p = 0{,}5$): $n = 50 \to \pm 13{,}86\text{ p.p.}$ ($\approx \pm 14\text{ p.p.}$); $n = 60 \to \pm 12{,}65\text{ p.p.}$ Os valores são indicativos, pois a fórmula supõe amostragem aleatória. Proporções de concordância (respostas 4 ou 5) serão reportadas com intervalo de Wilson de 95\%.
  \item \textbf{Regra de parada:} coleta de 10 dias corridos (12/11 a 22/11/2026); se houver $< 50$ respostas válidas, prorroga-se uma única vez, até 24/11/2026 às 12h (limite imposto pelo prazo de entrega dos resultados, 26/11 às 12h); persistindo, reporta-se a limitação.
\end{itemize}
